\ifdefined\AspersionSoloResultados
\else
\subsection{Acciones nominales de los chorros sobre el distribuidor}
\label{sec:man-asp-reacciones}

Para estimar las acciones hidráulicas se emplea la fuerza característica
de una boquilla a 4 bar, $F=4{,}04654$ N, calculada mediante cantidad de
movimiento. El líquido se acelera en la dirección del chorro; la acción
sobre el distribuidor tiene sentido contrario. Se considera el abanico
representado por una fuerza equivalente alineada con su eje. Esta
idealización permite estimar la escala de las cargas, pero no reemplaza
la integración de la cantidad de movimiento del abanico real.

Las cuatro boquillas superiores descargan hacia abajo. Su acción sobre
el distribuidor es ascendente. Las dos boquillas exteriores inclinadas
también aportan una componente ascendente. La suma nominal es
\begin{align}
 F_{z,nom}&=4F+2F\sin\alpha\notag\\
 &=4(4{,}04654)+2(4{,}04654)\sin43{,}637^\circ\notag\\
 &=\boxed{21{,}771\ \mathrm{N}}.
 \label{eq:man-asp-fz}
\end{align}
En cada ramal exterior, la magnitud de la resultante radial nominal es
\begin{align}
 F_{r,ext}&=4F+F\cos\alpha\notag\\
 &=4(4{,}04654)+4{,}04654\cos43{,}637^\circ
 =\boxed{19{,}115\ \mathrm{N}}.
\end{align}
Las resultantes de dos ramales diametralmente opuestos se compensan
globalmente solo si ambos conservan la misma orientación y descarga.
La compensación no elimina las fuerzas locales ni los momentos en sus
uniones. Para cinco boquillas interiores paralelas, la escala de la
resultante es
\begin{equation}
 F_{r,int}=5F=5(4{,}04654)=20{,}233\ \mathrm{N}.
\end{equation}
Si cambia su orientación, debe efectuarse la suma vectorial
$\boldsymbol{R}=-\sum_j\rho q_j\boldsymbol{v}_j$ y calcularse el momento
$\boldsymbol{M}_O=\sum_j\boldsymbol{r}_j\times\boldsymbol{R}_j$ respecto
del apoyo superior. Estas fuerzas no constituyen una envolvente
estructural: faltan el peso propio, el líquido contenido, las fuerzas
de presión en conexiones, el arranque, las mangueras y los casos de
ramales cerrados u obstruidos. Tampoco deben añadirse como reacciones
de rueda sin resolver previamente el camino de carga de la estructura.

\subsection{Cinemática nominal}
\label{sec:man-asp-cinematica}

Se adopta una velocidad de giro de 1,7 rev/min. La velocidad angular y
el período de una vuelta se obtienen convirtiendo a unidades SI:
\begin{align}
 \omega&=1{,}7\frac{\mathrm{rev}}{\mathrm{min}}
 \frac{2\pi\ \mathrm{rad}}{\mathrm{rev}}
 \frac{1\ \mathrm{min}}{60\ \mathrm{s}}
 =0{,}17802\ \mathrm{rad/s},\\
 T_g&=\frac{60}{1{,}7}=35{,}2941\ \mathrm{s}.
\end{align}
Como referencia, la velocidad tangencial de barrido sobre el diámetro
exterior es $u_D=\omega D/2=0{,}17802(0{,}710)=0{,}12640$ m/s, y sobre
el diámetro interior menor es $u_d=0{,}17802(0{,}450)=0{,}08011$ m/s.
Son velocidades geométricas del barrido, no del chorro. Para la mayor
distancia normal adoptada, el tiempo de vuelo rectilíneo estimado es
\begin{equation}
 t_v\simeq\frac{s_s}{v_j}
 =\frac{0{,}18485}{27{,}11354}
 =0{,}00682\ \mathrm{s}.
\end{equation}
Este tiempo es pequeño frente a los períodos mecánicos, lo que apoya
una primera evaluación cuasiestática de las huellas. La dispersión,
gravedad, interferencias y evolución real del abanico siguen requiriendo
una comprobación específica.
\fi

\ifdefined\AspersionSoloDiseno
\else
\section{Interpretación del caso base de diecinueve boquillas}
\label{sec:man-asp-interpretacion}

\subsection{Qué determina el cálculo hidráulico}

El balance permite definir una demanda hidráulica de referencia de
aproximadamente 174 L/min y 51 m de altura, con 19 boquillas simultáneas
y al menos 4 bar en sus tomas dentro del modelo simplificado. El caudal
indica cuánto líquido debe transportar el circuito; la altura expresa
la energía por unidad de peso que debe aportar la bomba. Ninguna de
estas dos magnitudes equivale a una eficacia de limpieza.

La alimentación DN40 presenta una velocidad de 2,02 m/s y reduce la
pérdida de la línea estacionaria a 2,64 m, frente a 20,85 m obtenidos
con DN25 para el mismo caudal y el trazado supuesto. Esta comparación
respalda DN40 como diámetro hidráulico preliminar. No selecciona por
sí sola el material, la clase de presión ni las conexiones.

La potencia al eje estimada es 2,58 kW. El consumo eléctrico de 3,03 kW
resulta de aplicar las eficiencias adoptadas, no de una curva de bomba
medida. La selección comercial deberá situar el punto de trabajo sobre
las curvas de caudal--altura, eficiencia y potencia, comprobar el
NPSH y verificar compatibilidad con el líquido a la temperatura prevista.

\subsection{Efecto de modificar la presión de boquilla}
\label{sec:man-asp-sensibilidad}

Manteniendo la misma boquilla y las propiedades del líquido,
$q\propto\sqrt{p_n}$ y $v_j\propto\sqrt{p_n}$. Por ello,
$F\simeq\rho qv_j\propto p_n$. Por ejemplo, pasar de 4 a 5 bar produce
\begin{align}
 \frac{q_5}{q_4}&=\sqrt{\frac54}=1{,}1180,
 &\frac{F_5}{F_4}&=\frac54=1{,}25.
\end{align}
El caudal individual aumenta aproximadamente 11,8\,\%, mientras la
escala de fuerza aumenta 25\,\%. La Tabla~\ref{tab:man-asp-sensibilidad}
expone la sustitución de distintas presiones en las mismas ecuaciones.

\begin{table}[htbp]
\centering\small
\caption{Sensibilidad de las boquillas a una presión uniforme impuesta.}
\label{tab:man-asp-sensibilidad}
\begin{tabular}{@{}rrrrr@{}}
\toprule
$p_g$ (bar) & $q$ (L/min) & $19q$ (L/min) & $v_j$ (m/s) & $F$ (N)\\\midrule
3 & 7{,}896 & 150{,}026 & 23{,}481 & 3{,}035\\
4 & 9{,}118 & 173{,}236 & 27{,}114 & 4{,}047\\
5 & 10{,}194 & 193{,}683 & 30{,}314 & 5{,}058\\
6 & 11{,}167 & 212{,}169 & 33{,}207 & 6{,}070\\
\bottomrule
\end{tabular}
\par\smallskip\footnotesize
Se considera la misma presión en las 19 boquillas y $C_v=0{,}95$.
No corresponde a una nueva resolución de la red ni a una curva de bomba.
\end{table}

Al aumentar la presión también aumenta el caudal de las tuberías y sus
pérdidas. No es correcto multiplicar toda la altura de bombeo por el
cuadrado de la razón de caudales: la elevación estática no depende de
$Q$, y la presión objetivo debe incorporarse como un término propio.
Para otra presión se requiere repetir el balance de cada tramo con su
nuevo Reynolds y con las pérdidas reales del filtro y la unión rotativa.
La mayor fuerza nominal tampoco demuestra una mejora proporcional de
la limpieza; esa relación queda fuera del alcance de este cálculo.

\subsection{Alcance geométrico y superficie de referencia}
\label{sec:man-asp-superficie}

Para interpretar el acceso se construye una envolvente geométrica
simplificada: superficie cilíndrica exterior, agujero interior y dos
caras anulares. Esta convención excluye el detalle de rodillos, jaula,
ranuras y superficies internas de los aros. Su área no es, por tanto,
el área real total que podría requerir limpieza.

El área lateral exterior se obtiene como perímetro por ancho axial:
\begin{equation}
 A_D=\pi DB=\pi(1{,}420)(0{,}515)=2{,}29745\ \mathrm{m^2}.
\end{equation}
Con una conicidad diametral de 1:30, el incremento de diámetro del
agujero a lo largo del ancho es $\Delta d=B/30=0{,}017167$ m. Su diámetro
medio es $d_m=d+\Delta d/2=0{,}908583$ m. Aproximando la longitud de la
generatriz por el ancho axial,
\begin{equation}
 A_d\simeq\pi d_mB=\pi(0{,}908583)(0{,}515)
 =1{,}47001\ \mathrm{m^2}.
\end{equation}
Para conservar una referencia simple se representan las dos caras con
el diámetro interior menor. El área convencional de una cara es
\begin{equation}
 A_c=\frac{\pi}{4}(D^2-d^2)
 =\frac{\pi}{4}(1{,}420^2-0{,}900^2)
 =0{,}94750\ \mathrm{m^2}.
\end{equation}
La suma de referencia resulta
\begin{equation}
 A_{ref}=A_D+A_d+2A_c
 =2{,}29745+1{,}47001+2(0{,}94750)
 =\boxed{5{,}66247\ \mathrm{m^2}}.
 \label{eq:man-asp-aref}
\end{equation}
Se trata de una aproximación explícita. Para un inventario superficial
de precisión deben emplearse las dos caras con sus diámetros reales y
la generatriz cónica, además del resto de la geometría del rodamiento.

En un escenario ideal, si todo el exterior, todo el agujero y la cara
superior fueran alcanzables, y se descontaran tres zonas laterales de
retención representadas por rectángulos de $80\times30$ mm, el área sería
\begin{align}
 A_{ret}&=3(0{,}080)(0{,}030)=0{,}00720\ \mathrm{m^2},\\
 A_{pot}&=A_D-A_{ret}+A_d+A_c\notag\\
 &=2{,}29745-0{,}00720+1{,}47001+0{,}94750
 =4{,}70777\ \mathrm{m^2},\\
 \lambda_{pot}&=\frac{A_{pot}}{A_{ref}}
 =\frac{4{,}70777}{5{,}66247}=0{,}8314.
 \label{eq:man-asp-lambda}
\end{align}
El 83,14\,\% resultante es una fracción potencial de esa envolvente
convencional, bajo las condiciones indicadas. No es una cobertura
demostrada: la franja nominal de 13,76 mm identificada en la
Ecuación~\eqref{eq:man-asp-hueco}, las sombras y las trayectorias pueden
reducirla. La cara inferior no se incluye en $A_{pot}$, porque el caso
base no dispone una aplicación ascendente acreditada hacia ella.

Como referencia para estudiar posteriormente dicha cara, tres sectores
de apoyo de radios 638 y 688 mm y ángulo de 24$^\circ$ tendrían un área
proyectada conjunta
\begin{equation}
 A_{\mathrm{ap}}=3\pi(0{,}688^2-0{,}638^2)\frac{24}{360}
 =0{,}04166\ \mathrm{m^2}.
\end{equation}
Esta proyección representa una hipótesis geométrica para descontar zonas
ocupadas; no determina por sí sola el contacto resistente real entre
el rodamiento y el liner. Tampoco acredita que la superficie inferior
restante sea accesible desde las boquillas existentes.

La cobertura geométrica y la remoción deben mantenerse como indicadores
independientes. Una superficie visible puede recibir un chorro
insuficiente, y una superficie inicialmente sombreada puede requerir
otra orientación o una operación complementaria. El porcentaje de
limpieza solo podrá establecerse después de definir un criterio de
aceptación y medir el contaminante residual.

\subsection{Caudal recirculado y volumen de líquido}

Un caudal de 173,99 L/min durante diez minutos representa un volumen
acumulado que ha circulado por las boquillas de
\begin{equation}
 V_{circ}=Qt=(173{,}99168)(10)=1739{,}92\ \mathrm{L}.
\end{equation}
En un circuito recirculado ese volumen puede contener varias pasadas
del mismo líquido. Por tanto, no exige un tanque de 1740 L ni equivale
al consumo de agua nueva. El volumen útil del depósito debe permitir
mantener el nivel de succión, alojar el inventario de tuberías y retorno,
separar aire y contaminantes y absorber las variaciones durante el ciclo.
La reposición depende de arrastre, evaporación, purgas y pérdidas.

\subsection{Síntesis y condiciones de aplicación}
\label{sec:man-asp-conclusion}

La Tabla~\ref{tab:man-asp-sintesis} ordena los resultados principales del
caso base y su significado dentro del diseño.

\begin{table}[htbp]
\centering\small
\caption{Resultados principales y significado dentro del diseño.}
\label{tab:man-asp-sintesis}
\begin{tabularx}{\linewidth}{@{}>{\raggedright\arraybackslash}p{3.5cm} >{\raggedright\arraybackslash}p{2.6cm} >{\raggedright\arraybackslash}X@{}}
\toprule Magnitud & Resultado & Interpretación\\\midrule
Cantidad simultánea & 19 boquillas & Configuración del caso base hidráulico; cobertura por completar.\\
Caudal de alimentación & 173,99 L/min & Resultado de continuidad y presión en la red simplificada.\\
Presión mínima de toma & 4,00 bar & Condición impuesta para el predimensionamiento.\\
Altura de bombeo & 50,75 m & Incluye presión, desnivel, pérdidas y carga cinética.\\
Alimentación principal & DN40; 2,02 m/s & Alternativa hidráulica preliminar frente a DN25.\\
Potencia al eje & 2,58 kW & Demanda con eficiencia de bomba adoptada de 0,55.\\
Consumo de bombeo & 3,03 kW & Estimación con eficiencia de motor de 0,85.\\
NPSH disponible & 6,01 m & Escenario a 2335 m, 62 $^\circ$C y nivel inundado de 0,50 m.\\
Franja axial nominal & 13,76 mm & Discontinuidad geométrica; no se declara cobertura completa.\\
\bottomrule
\end{tabularx}
\end{table}

El desarrollo establece una base cuantitativa para especificar la bomba,
contrastar diámetros y preparar la verificación de cobertura. Los
decimales de las operaciones permiten seguir la sustitución numérica;
no representan precisión experimental. Para especificación preliminar
se retienen aproximadamente 174 L/min y 51 m.

El cierre del dimensionamiento requiere confirmar las propiedades de
la solución de limpieza; las curvas de boquilla, filtro, unión rotativa
y bomba; el trazado y los accesorios; y el nivel mínimo del depósito.
Con esos datos debe refinarse la distribución de presión incorporando
los cambios de velocidad y las conexiones de las tomas. El cierre de
cobertura requiere corregir la franja lateral y comprobar interferencias,
sombras y acceso inferior; la eficacia de lavado queda fuera
del alcance del caso analítico. Estas verificaciones delimitan la aplicación del
caso analítico y no invalidan las operaciones de predimensionamiento
desarrolladas.
\fi